\chapter*{Conclusion}
\addcontentsline{toc}{chapter}{Conclusion}
\label{chap:conclusion}

This thesis addressed the challenge of classifying causal variable roles in observational data, focusing on the ADIA Lab Causal Discovery Challenge where the goal is to assign one of eight labels (Confounder, Mediator, Collider, Independent, Cause of X/Y, Consequence of X/Y) to variables given their relationships with designated treatment $X$ and outcome $Y$. The problem is difficult because several role pairs share nearly identical statistical signatures -- for example, Confounder and Mediator both produce conditional independence $X \perp Y \mid K$ -- requiring models to distinguish subtle structural patterns rather than merely testing for dependence.

The proposed solution is a \textbf{dual-pipeline neural architecture} combining 1D edge tensor processing with 2D scatter density visualization, achieving \textbf{81.082\%} Balanced Accuracy on the private test set. This represents a \textbf{+4.38\% improvement} over the competition first-place solution (76.70\%) and a \textbf{+7.12\% improvement} over a faithful reimplementation baseline (73.96\%). The model makes five specific technical contributions, each validated by systematic ablation: (1)~8-channel edge tensors with multi-bandwidth kernel regression and ANM residuals (+1.89\%); (2)~structural attention bias encoding six topological relationship types (+2.65\%); (3)~2D scatter density pipeline capturing joint distributional patterns (+0.88\%); (4)~$\lambda$-weighted dual-loss training ensuring edge context quality (+0.59\%); and (5)~X/Y remap data augmentation multiplying effective training data 11-fold (+1.55\%).

Beyond the numerical improvements, this thesis establishes four methodological insights with broader relevance to causal discovery systems. \textbf{First}, scalar feature compression is irreversible -- once the full shape of conditional relationships is reduced to correlation coefficients or independence test statistics, the information distinguishing Confounder from Mediator is permanently lost. The 8-channel edge tensor preserves the complete functional curve, enabling Conv1D to learn discriminative patterns invisible to scalar methods. \textbf{Second}, structured inductive bias consistently outperforms unconstrained capacity. The structural attention bias requires only six additional parameters yet provides +2.65\% gain, while failed experiments adding thousands of unconstrained parameters (parallel transformers, node-centric attention, deep sets) produced negative gains of -1.19\% to -1.90\%. \textbf{Third}, 1D curve representations and 2D density images are complementary. Edge tensors excel at directional asymmetry (Collider detection), while scatter density images capture joint patterns (Confounder/Mediator distinction); neither representation subsumes the other. \textbf{Fourth}, edge context is a structural necessity, not an auxiliary regularization term. Removing the edge pipeline causes performance to collapse from 80.47\% to 51.51\% -- below even classical machine learning baselines -- confirming that the binary edge loss forces embeddings to encode causally meaningful graph structure essential for node classification.

\textbf{Limitations and future work.} Despite achieving state-of-the-art performance, three challenges remain. Mediator (73.71\%) and Collider (78.70\%) remain the hardest classes with substantial room for improvement. The model has been validated only on synthetic data from the ADIA Lab competition; real-world causal graphs may exhibit different noise characteristics, missing data patterns, or temporal dynamics not present in the training distribution. Finally, the architecture requires GPU resources for both training and inference, limiting deployment in resource-constrained environments. Future work should validate the approach on real-world causal inference benchmarks, extend the method to time series causal discovery where temporal ordering provides additional constraints, and investigate distillation techniques to reduce model size while preserving accuracy.

This work demonstrates that deep learning architectures can advance causal discovery beyond the limitations of classical methods when design choices are grounded in the structural properties of causal graphs. The 81.082\% Balanced Accuracy establishes a new state-of-the-art for the ADIA Lab Challenge and provides a foundation for automated causal analysis systems that can reason about variable roles in complex systems.
